% Generated draft: requires mathematical and source review.
\documentclass[11pt]{article}
\usepackage[margin=25mm]{geometry}
\usepackage{fontspec}
\setmainfont{Latin Modern Roman}
\usepackage{amsmath,amssymb,mathtools,unicode-math}
\setmathfont{Latin Modern Math}
\usepackage{enumitem,needspace,xurl,hyperref}
\setlength{\parindent}{0pt}
\setlength{\parskip}{5pt}
\setlength{\emergencystretch}{3em}
\newcommand{\N}{\mathbb{N}}
\newcommand{\Z}{\mathbb{Z}}
\newcommand{\Q}{\mathbb{Q}}
\newcommand{\R}{\mathbb{R}}
\newcommand{\C}{\mathbb{C}}
\newcommand{\F}{\mathbb{F}}
\newcommand{\A}{\mathbb{A}}
\newcommand{\PP}{\mathbb{P}}
\newcommand{\E}{\mathbb{E}}
\newcommand{\eps}{\varepsilon}
\newcommand{\dd}{\,\mathrm{d}}
\begin{document}
\section*{UnsolvedMath ID: 30003488}
\textbf{Source problem code:} catalogue entry
\subsection{Definitions and mathematical statement}
Let \(X\) be a finite set of species labels, and let \(\sigma=(\tau,\ell)\) denote a rooted species tree with topology \(\tau\) and branch-length parameters \(\ell\), under the multispecies coalescent model. The supplied record does not specify the admissible tree class, parameter domain, branch-length units, or sampling convention. These must be fixed before the questions below define a single mathematical problem; in particular, no assumption about strictly positive versus zero branch lengths is imposed here.

For a fixed sampling convention, let \(L\) be the set of sampled gene-lineage labels and let \(G_\sigma\) be the random gene tree generated by the model. A split of \(L\) is an unordered bipartition \(A\mid(L\setminus A)\), where \(\varnothing\ne A\subsetneq L\). Say that the split is displayed when deleting an edge of the underlying unrooted gene tree separates its leaves into those two parts. For the collection \(\mathcal S\) of observed splits, define
\[
 p_s(\sigma)=\Pr_\sigma\{G_\sigma\text{ displays }s\},\qquad
 p(\sigma)=(p_s(\sigma))_{s\in\mathcal S}.
\]
The sampling convention and the collection \(\mathcal S\), including whether trivial splits are included, are not specified by the record.

\begin{enumerate}
\item \textbf{Branch-length identifiability.} Determine whether the branch lengths of \(\sigma\) are identifiable from \(p(\sigma)\). For example, after fixing a topology \(\tau\), an admissible parameter domain \(\Theta_\tau\), and a set \(E_*\) of branches whose lengths are to be recovered, the global, fixed-topology interpretation is the implication
\[
 \forall\ell,\ell'\in\Theta_\tau,\qquad
 p(\tau,\ell)=p(\tau,\ell')
 \ \Longrightarrow\ 
 \ell|_{E_*}=\ell'|_{E_*}.
\]
This is a conditional formulation, not a choice of missing hypotheses: the record does not say whether topology is known, which branches are included, or whether global, generic, or local identifiability is intended. If topology is unknown, a corresponding formulation must also specify how branches of two candidate trees are matched.
\item \textbf{Root estimation.} Determine whether an improved NJst algorithm can estimate the root of the species tree. The source does not define NJst, the permitted improvement, its input data, the meaning of the root parameter, or the required estimation guarantee. In particular, it does not specify whether the target is a root attachment position on an unrooted topology or also a metric location, nor whether estimation means exact recovery from population quantities or statistical consistency from finite samples. A precise algorithmic assertion therefore remains unresolved.
\item \textbf{Higher-degree phylogenetic invariants.} For a fixed topology \(\tau\), a polynomial split-probability invariant may be expressed as a nonzero polynomial
\[
 f\in\R[z_s:s\in\mathcal S]
 \quad\text{such that}\quad
 f\bigl(p(\tau,\ell)\bigr)=0
 \quad\text{for every }\ell\in\Theta_\tau.
\]
Ask whether higher-degree phylogenetic invariants exist for the intended model. The record does not specify the probability coordinates, coefficient field, reference degree, or whether the intended invariants must provide relations beyond previously known ones. Thus this displayed definition supplies a conditional interpretation in split coordinates; it does not resolve the intended meaning of ``higher-degree.''
\end{enumerate}

The supplied dated assessment labels the combined problem partially solved: it reports topology identifiability from split probabilities in the multispecies coalescent setting, branch-length recovery only in restricted or local cases, and broader branch-length, root, and invariant questions as remaining open. These are supplied curation claims, not independently verified results, and the record supplies no precise hypotheses for the reported results.
\subsection{Short English statement}
The problem asks whether split probabilities under the multispecies coalescent determine species-tree branch lengths, whether an improved NJst algorithm can estimate the root, and whether higher-degree phylogenetic invariants exist. Sampling, parameter domains, recoverable branches, estimation guarantees, and the intended invariant degree remain unspecified. The supplied partially solved status requires independent verification and human mathematical review.
\subsection{Sources}
\begin{itemize}
\item Supplied local UnsolvedMath registry record, permanent ID \(30003488\), title ``Branch Length and Root Identifiability from Split Probabilities,'' source problem code \texttt{catalogue entry}, in the Oberwolfach Reports Open Problems collection. The record gives proposed year 2017 and names no individual proposer.
\item Supplied source citation: \emph{Algebraic Statistics (2018)}, Oberwolfach Reports, DOI 10.4171/owr/2017/20, pp.~1207--1279. \url{https://doi.org/10.4171/owr/2017/20}. The background describes these reports as documenting mathematical workshops and their open-problem discussions. The supplied literature-source label for this same link is ``Oberwolfach Reports 2018.''
\item Supplied literature source: ``Split probabilities and species tree inference.'' \url{https://pmc.ncbi.nlm.nih.gov/articles/PMC5766424/}.
\item Supplied literature source: ``Chor and Steel, tree split probabilities.'' \url{https://doi.org/10.1016/j.jtbi.2015.03.024}.
\item Literature assessment date: \texttt{2026-08-21}. This is supplied dated curation, not independently verified. Its assessment reports proved topology identifiability, restricted or local branch-length recovery, and broader unresolved branch-length, root, and invariant questions. The background explicitly calls this dated literature triage derived from the supplied OWR curation and requests independent verification of time-sensitive status and cited claims.
\item The registry labels the statement \texttt{corrected\_verified}, explaining that mathematical content was preserved while notation, typography, and sentence structure were repaired for publication. This supplied editorial label does not establish mathematical correctness or independent source verification.
\end{itemize}
\end{document}
